\section{Goods per Agent}\label{sec:goods}
%% FUNCTION: Shows how the EFX rate of round robin changes with the number of goods per agent r
%% when the number of goods is divisible by the number of agents (q = 0; m = 8, 12, 16 for n = 4).
%% More goods per agent help unless agents agree strongly. Special case: if every agent receives
%% at most two goods and all favourite goods differ, round robin is guaranteed to be EFX (proven);
%% the guarantee breaks at m = 2n + 1. Since this break coincides with the first remainder, the
%% section leads to divisibility.

The number of goods per agent says how many rounds the agents pick. Two flatmates who split four
pieces of furniture pick twice each; four heirs who split sixteen items pick four times each. In
our design, it is the number of full rounds $r=\lfloor m/n\rfloor$. To separate it from
divisibility, we keep the number of goods divisible by the number of agents ($q=0$) and vary $r$
from $2$ to $4$, which means $m=8$, $12$ and $16$ goods for $n=4$ agents. Agreement and dispersion
stay at the reference levels, Spearman $0.48$ ($\rho=0.5$) and Gini $0.52$ ($\theta=1$). The case
$r=1$ occurs in our design only with $q>0$; Section~\ref{sec:bundlesize} covers it.

\subsection{EFX Rate}\label{sec:goods-rate}

As Figure~\ref{fig:goods} shows, Round Robin achieves EFX more often when every agent receives
more goods. With two goods per agent, the EFX rate is $0.212$, with $15{,}755$ failures among
$R=20{,}000$ instances. It rises to $0.293$ with three goods and to $0.409$ with four goods, where
$11{,}828$ instances fail (measured).

\begin{figure}[ht]
  \centering
  \includegraphics{fig43_goods_rate.pdf}
  \caption{EFX rate of Round Robin against the number of goods per agent $r$, for $n=4$ agents
  with $m=4r$ goods ($q=0$), Spearman $0.48$ ($\rho=0.5$) and Gini $0.52$ ($\theta=1$). Each
  point rests on $R=20{,}000$ instances.}
  \label{fig:goods}
\end{figure}

This direction does not hold everywhere. We fixed a counting rule before counting: for each
number of agents, level of agreement and level of dispersion, we follow the EFX rate over
$r=2,3,4$ with $q=0$, and a course counts as rising if the rate increases strictly at both steps.
Because we knew individual cells beforehand, we treat the count as exploratory. Of the $120$
courses, $20$ lie at exactly $1$ throughout; these are all courses at Gini $0.06$ ($\theta=0.1$),
as Proposition~\ref{prop:theta-limit} predicts. Of the $80$ remaining courses below full
agreement, $70$ rise. The exceptions concentrate where agreement is strong: up to Spearman $0.48$
($\rho=0.5$), $59$ of $60$ courses rise, at Spearman $0.73$ ($\rho=0.75$) only $11$ of $20$. In
these exceptions, the EFX rate dips slightly at $r=3$ and recovers at $r=4$. Under full
agreement, the direction reverses: none of the $20$ courses rises, and the EFX rate falls with
every additional round. For $n=4$ and Gini $0.28$ ($\theta=0.5$), for instance, it falls from
$0.217$ with two goods per agent to $0.047$ and $0.010$. We have no explanation for this
reversal. In summary, more goods per agent help Round Robin as long as the agents do not agree
too strongly (measured).

\subsection{At Most Two Goods and Distinct Favourite Goods}\label{sec:bundlesize}

With few goods per agent, there is a case in which Round Robin is guaranteed to achieve EFX. It
depends on whether the agents have different favourite goods. We call an instance
\textit{collision-free} if all agents have different favourite goods. For $\rho=1$, all agents
have the same favourite good, so collision-freeness never occurs.

\begin{lemma}\label{lem:bundlesize}
If an instance is collision-free, then under Round Robin an agent $i$ can violate EFX towards an
agent $j$ only if $|A_j|\geq3$.
\end{lemma}

\begin{proof}
Without a collision, every agent receives their favourite good in the first round, so
$v_i(A_i)\geq v_i(g)$ for every good $g$. If $|A_j|=1$, EFX holds trivially for this pair. If
$|A_j|=2$, removing either good from $A_j$ leaves a single good $g$, and
$v_i(g)\leq v_i(A_i)$.
\end{proof}

\begin{corollary}\label{cor:bundlesize}
If an instance is collision-free and $m\leq2n$, the Round Robin allocation is EFX.
\end{corollary}

\begin{proof}
For $m\leq2n$, Round Robin gives every agent at most two goods, so Lemma~\ref{lem:bundlesize}
applies to every pair.
\end{proof}

We checked the corollary in every stored cell: in all $336$ cells with $m\leq2n$ and $\rho<1$,
over all six levels of dispersion, each of the $1{,}388{,}388$ collision-free instances is EFX.
The proof uses only the favourite goods and the bundle sizes, so the guarantee does not depend on
our distribution of values. To our knowledge, it has not been stated explicitly before.

A related observation explains why the EFX rate at $m=n+1$ does not depend on dispersion
(Section~\ref{sec:dispersion-rate}). With one good more than agents, only one agent receives two
goods. For every other pair, EFX holds trivially or requires $v_i(A_i)\geq\max_{g\in A_j}v_i(g)$,
a comparison of single goods that depends only on the rankings. Since dispersion does not change
rankings, the EFX rate at $m=n+1$ is the same for every level of dispersion (proven; in the
stored cells it is identical instance by instance).

The condition of the corollary quickly becomes rare. Table~\ref{tab:collisionfree} shows the
share of collision-free instances at $m=2n$, which does not depend on dispersion. It falls from
$0.749$ for $n=2$ and Spearman $0$ to $0.010$ for $n=5$ and Spearman $0.73$, which are only $193$
of $20{,}000$ instances.

\begin{table}[ht]
  \centering
  \caption{Share of collision-free instances at $m=2n$ under Round Robin, $R=20{,}000$ per cell.}
  \label{tab:collisionfree}
  \begin{tabular}{lcccc}
    \toprule
    & \multicolumn{4}{c}{Spearman ($\rho$)} \\
    $n$ & $0$ ($0$) & $0.24$ ($0.25$) & $0.48$ ($0.5$) & $0.73$ ($0.75$) \\
    \midrule
    2 & 0.749 & 0.650 & 0.535 & 0.379 \\
    3 & 0.553 & 0.398 & 0.256 & 0.126 \\
    4 & 0.405 & 0.231 & 0.115 & 0.038 \\
    5 & 0.299 & 0.140 & 0.051 & 0.010 \\
    \bottomrule
  \end{tabular}
\end{table}

Figure~\ref{fig:goods-break} shows what happens beyond $m=2n$ for $n=4$ agents, restricted to
collision-free instances. Up to $m=8$, the EFX rate is exactly $1$, as the corollary states. As
soon as one agent receives three goods, it drops: at $m=9$, to $0.830$ for Gini $0.52$
($\theta=1$), with $467$ failures among $2{,}742$ collision-free instances, and to $0.682$ for
Gini $0.06$ ($\theta=0.1$), with $871$ failures. By Lemma~\ref{lem:bundlesize}, the only pairs at
risk are those in which an agent with two goods compares their bundle with a bundle of three
goods. To the right of the break, the EFX rate recovers unevenly, to $0.941$ at $m=16$ for Gini
$0.52$, and at Gini $0.06$ every divisible number of goods returns to exactly $1$, as in
Section~\ref{sec:dispersion-limit}. The break also deepens with agreement: among collision-free
instances at $m=9$ and Gini $0.52$, the EFX rate is $0.948$, $0.912$, $0.830$ and $0.696$ for
Spearman $0$, $0.24$, $0.48$ and $0.73$ (measured). Over all instances, the break is hardly
visible, because most instances contain a collision: the EFX rate only falls from $0.212$ at
$m=8$ to $0.179$ at $m=9$.

\begin{figure}[ht]
  \centering
  \includegraphics{fig43_goods_break.pdf}
  \caption{EFX rate of Round Robin among collision-free instances against the number of goods
  $m$, for $n=4$ agents and Spearman $0.48$ ($\rho=0.5$), at Gini $0.52$ ($\theta=1$) and Gini
  $0.06$ ($\theta=0.1$). Left of the dashed line, no agent receives more than two goods. The
  number of collision-free instances per point ranges from $861$ at $m=5$ to $4{,}928$ at $m=16$
  out of $R=20{,}000$.}
  \label{fig:goods-break}
\end{figure}

At $m=2n+1$, two things change at once: the first bundle with three goods appears, and the
remainder changes from $q=0$ to $q=1$. Among collision-free instances, the break is a pure effect
of bundle size, because for $m\leq2n$ many numbers of goods with $q>0$ occur and all of them
reach exactly $1$. Over all instances, both effects mix; Section~\ref{sec:divisibility} isolates
divisibility.

\subsection{Severity of Failures}\label{sec:goods-severity}

Figure~\ref{fig:goods-survival} shows the survival curves for two, three and four goods per
agent. A curve for more goods per agent never lies below a curve for fewer goods, so more rounds
reduce both the number and the depth of failures. At $\alpha=0.6$, the share of instances that
reach at least this level is $0.468$, $0.651$ and $0.787$. The conditional $5\%$ quantile rises
from $0.14$ to $0.23$ and $0.30$, and the median of $\astar$ among the instances that are not EFX
rises from $0.47$ to $0.60$ and $0.69$. Severe failures become rare: the share of instances with
$\astar<0.2$ falls from $0.090$ with two goods per agent to $0.024$ and $0.008$ (measured).

\begin{figure}[ht]
  \centering
  \includegraphics{fig43_goods_survival.pdf}
  \caption{Share of instances with $\astar\geq\alpha$ under Round Robin, one curve per number of
  goods per agent, for $n=4$ agents with $m=4r$ goods ($q=0$), Spearman $0.48$ ($\rho=0.5$) and
  Gini $0.52$ ($\theta=1$). Each curve rests on $R=20{,}000$ instances. The point on each curve
  marks the $5\%$ quantile of $\astar$ among the instances that are not EFX.}
  \label{fig:goods-survival}
\end{figure}

\paragraph{Relation to literature.}
Our results fit the asymptotic picture for independent values. \textcite[Thm.~3.1]{manurangsi2021}
show that Round Robin even achieves envy-freeness with high probability once the number of goods
is of order $n\log n/\log\log n$, for independent values whose density on $[0,1]$ is bounded.
Their result does not formally apply to our lognormal values, and it concerns many agents, but it
points in the same direction as our measurements at low agreement: more goods per agent help.
\textcite[p.~17]{manurangsi2021} also argue that standard Round Robin with one round and a large
remainder fails EFX with constant probability, because the last agent may lose their favourite
good to one of the agents who pick an extra good. This is the collision mechanism behind
Lemma~\ref{lem:bundlesize}, in asymptotic form. The existence of an EFX allocation is proven in
most cells with $m\leq2n$: for $n\leq3$ \parencite{plaut2020,chaudhury2024}, for $m\leq n+2$
\parencite{amanatidis2020multiple} and for $m\leq n+3$ \parencite{mahara2024general}. For $n=4$,
it is open from $m=8$ on, so all three cells of Figure~\ref{fig:goods} lie in the region where we
assume existence (Section~\ref{sec:existence}).

\paragraph{Recommendation.}
If every person receives at most two goods and all favourite goods differ, Round Robin delivers
EFX for certain, and both conditions can be checked before picking starts; beyond that, more goods
per person make EFX more likely, unless everyone agrees on what is valuable. Because the
guarantee breaks exactly where the first agent receives one good more than the others, the next
section asks what divisibility does.
